% Generated draft: requires mathematical and source review.
\documentclass[11pt]{article}
\usepackage[margin=25mm]{geometry}
\usepackage{fontspec}
\setmainfont{Latin Modern Roman}
\usepackage{amsmath,amssymb,mathtools,unicode-math}
\setmathfont{Latin Modern Math}
\usepackage{enumitem,needspace,xurl,hyperref}
\setlength{\parindent}{0pt}
\setlength{\parskip}{5pt}
\setlength{\emergencystretch}{3em}
\newcommand{\N}{\mathbb{N}}
\newcommand{\Z}{\mathbb{Z}}
\newcommand{\Q}{\mathbb{Q}}
\newcommand{\R}{\mathbb{R}}
\newcommand{\C}{\mathbb{C}}
\newcommand{\F}{\mathbb{F}}
\newcommand{\A}{\mathbb{A}}
\newcommand{\PP}{\mathbb{P}}
\newcommand{\E}{\mathbb{E}}
\newcommand{\eps}{\varepsilon}
\newcommand{\dd}{\,\mathrm{d}}
\begin{document}
\section*{UnsolvedMath ID: 2100708}
\textbf{Source problem code:} catalogue entry
\subsection{Definitions and mathematical statement}
Consider a semiclassical quantum integrable family $T_{1,\hbar},\ldots,T_{n,\hbar}$ with commuting self-adjoint operators and classical principal symbols $f_1,\ldots,f_n$ on a symplectic $2n$-manifold. The usual classical integrability condition is $\{f_i,f_j\}=0$, with independent differentials on the regular locus; put $F=(f_1,\ldots,f_n)$. In a region where the joint spectrum is discrete, its points are tuples $\lambda=(\lambda_1,\ldots,\lambda_n)$ for which a nonzero common eigenvector satisfies
\[
T_{j,\hbar}u=\lambda_j u\qquad(1\le j\le n).
\]
This is the standard discrete semiclassical interpretation, not a source-specified choice between pseudodifferential, Toeplitz, or other quantizations.

A classical point $p$ is singular if $r=\operatorname{rank}\dd F(p)<n$. In the usual Williamson convention, it is nondegenerate when the transverse quadratic Hamiltonians on the $2(n-r)$-dimensional symplectic normal space span the required Cartan subalgebra of the symplectic Lie algebra; otherwise it is degenerate. This concerns the classical integrable singularity, rather than multiplicity or degeneracy of a quantum eigenvalue. The source does not supply its exact nondegeneracy convention.

Pelayo's question is whether one can detect degeneracy of a singularity from the joint spectrum alone. A precise version seeks a spectral criterion, based on specified joint spectral data as $\hbar\to0$ near the relevant classical value $c=F(p)$, that decides the intended classical degeneracy predicate. The extract does not say whether it concerns a specified singular point, the presence of at least one degenerate point in a singular fiber, or the global presence of degenerate singularities. It also does not specify whether one knows exact spectra, multiplicities, complete asymptotic data, or data only modulo a stated semiclassical error. These choices must be fixed before a universal recoverability assertion can be quantified.

The supplied summary reports nondegenerate spectral distinctions in some settings, including elliptic, hyperbolic, and focus--focus patterns, and describes degenerate spectral detection as unresolved by its check. It mentions differing clustering rates, but provides no hypotheses or theorem making such a rate a universal degeneracy test. These are attributed partial-progress and source-verification claims, not a criterion proved here.
\subsection{Short English statement}
Can semiclassical joint eigenvalue data reveal that an underlying classical integrable singularity is degenerate? The question is about classical nondegeneracy, not repeated eigenvalues. The source leaves the quantization, spectral accuracy, locality, and point-versus-fiber target unspecified, and reports only partial nondegenerate results without a verified general degeneracy criterion.
\subsection{Sources}
\begin{itemize}
\item Local source registry: \url{data/unsolvedmath/problems.json}, permanent UnsolvedMath ID $2100708$, source code catalogue entry, including the supplied PARTIAL-PROGRESS summary. Dataset attribution: Ulam AI and UnsolvedMath Contributors; original-source attribution and rights are retained.
\item \emph{Open Problems, Questions, and Challenges in Finite-Dimensional Integrable Systems} (2018), Problem 7.8, attributed to \'{A}. Pelayo: \url{https://arxiv.org/abs/1804.03737}. The metadata reports source-TeX extraction and July 29, 2026 access, with status and rights review flags; no source-paper author list is supplied in this record.
\item The supplied summary refers generically to semiclassical spectral clustering results for nondegenerate singularities, but provides no precise additional reference, hypotheses, or verified criterion for the degenerate case. No linked source was inspected for this draft.
\end{itemize}
\end{document}
